% Chapter 5

\chapter{面向越狱攻击的分层防御方案}

\section{防御问题定义与设计目标}

第四章的静态实验结果表明，大语言模型的安全脆弱性并不集中于单一提示模板，而是同时体现在语义包装、结构重构、上下文诱导与输出边界模糊等多个层面。特别是认知误导类与指令重构类攻击在多个模型上重复暴露出较强突破能力，而大量不确定样本又说明当前安全对齐并未形成足够稳定的硬边界。由此可见，越狱防御问题不能被简单理解为“拦截敏感词”或“增加拒答模板”，而应被视为一个覆盖输入、交互与输出全过程的系统性风险控制问题\cite{rebedea2023nemo,phute2024llmselfdefense,li2023textpurification}。

基于这一认识，本文提出面向越狱攻击的分层防御方案，其总体目标是在尽可能保留模型正常可用性的前提下，对不同阶段出现的风险信号进行分层识别、逐级干预与过程留痕。与单点式防御不同，该方案并不期望某一机制独立解决全部安全问题，而是通过输入层、交互层与输出层的协同作用，使风险能够在进入前、推进中和生成后三个阶段分别被识别、削弱与兜底。

从设计原则上看，本文防御方案围绕以下四个目标展开：第一，提升对显式高风险输入与典型越狱模式的早期识别能力，在模型调用前尽可能阻断明显恶意请求；第二，对多轮诱导、规则豁免、角色操控和持续性试探等过程性风险进行动态监测，避免攻击在连续交互中逐步放大；第三，对已经生成的潜在危险内容实施输出侧修正或替换，降低前置防线失效后的暴露风险；第四，为后续误报复核、规则优化与防御评估保留审计信息，使防御过程具备可解释性与可迭代性。

\begin{figure}[htbp]
\centering
\resizebox{0.98\textwidth}{!}{\begin{tikzpicture}[
    x=1pt,y=1pt,font=\small,>=stealth,
    block/.style={draw=#1!75!black, fill=#1!12, rounded corners=9pt, line width=0.9pt, align=center},
    flow/.style={->, line width=0.95pt, draw=gray!75}
]
\node[block=orange, minimum width=110pt, minimum height=38pt] (in) at (100,102) {\textbf{输入前}\\显式恶意意图\\提示注入信号};
\node[block=teal, minimum width=110pt, minimum height=38pt] (mid) at (294,102) {\textbf{交互中}\\多轮诱导\\上下文累积};
\node[block=red, minimum width=110pt, minimum height=38pt] (out) at (488,102) {\textbf{输出后}\\危险步骤\\边界性内容泄露};
\draw[flow] (in.east) -- node[above, font=\scriptsize] {风险推进} (mid.west);
\draw[flow] (mid.east) -- node[above, font=\scriptsize] {继续放大} (out.west);
\node[anchor=west, font=\bfseries\small] at (18,164) {越狱防御问题定义图};
\node[anchor=west, font=\scriptsize, text=gray!70!black] at (18,28) {风险并非只出现在单一点位, 而是沿输入、交互、输出全过程传播};
\end{tikzpicture}
}
\caption{越狱防御问题定义图}
\label{fig:chapter5-defense-problem}
\end{figure}

\section{分层防御总体框架}

\subsection{分层防御的总体思路}

本文防御方案采用“输入层防御 + 交互层防御 + 输出层防御”的三层协同结构，其核心思想是在模型调用链的不同阶段分别承担不同职责。输入层位于模型调用之前，主要负责识别显式危险意图、提示注入信号与典型越狱模式，是整个系统的第一道防线；交互层位于过程控制位置，主要负责跟踪风险累积、诱导行为与轮次推进中的状态变化，是针对动态对齐失效的过程控制层；输出层位于模型生成之后，主要负责复核回复中的危险内容、实施脱敏替换并记录审计结果，是最后一道内容安全兜底层。

这一结构的必要性在于，不同攻击路径作用于模型的方式并不相同。角色操控、规则豁免与危险意图的显式表达，更适合在输入阶段进行识别；示例诱导、上下文累积与连续试探，则需要在交互过程中持续跟踪；而一旦模型已经生成潜在危险内容，则必须依赖输出层进行再检测与再处理。若仅依赖单层防御，不仅容易出现遗漏，也会使防御动作过于粗糙，从而在安全性与可用性之间形成更大冲突\cite{rebedea2023nemo,phute2024llmselfdefense}。

\begin{figure}[htbp]
\centering
\resizebox{0.98\textwidth}{!}{\begin{tikzpicture}[
    x=1pt,y=1pt,font=\small,>=stealth,
    panel/.style={draw=gray!35, rounded corners=14pt, fill=gray!4, line width=1pt},
    block/.style={draw=#1!75!black, fill=#1!12, rounded corners=9pt, line width=0.9pt, align=center},
    flow/.style={->, line width=0.95pt, draw=gray!75}
]
\draw[panel] (20,34) rectangle (716,210);
\node[block=gray, minimum width=92pt, minimum height=34pt] (user) at (78,122) {\textbf{用户输入}};
\node[block=blue, minimum width=118pt, minimum height=54pt] (in) at (218,122) {\textbf{输入层防御}\\规则匹配\\分类打分};
\node[block=teal, minimum width=118pt, minimum height=54pt] (mid) at (388,122) {\textbf{交互层防御}\\状态跟踪\\风险累积};
\node[block=red, minimum width=118pt, minimum height=54pt] (out) at (558,122) {\textbf{输出层防御}\\内容复核\\替换与脱敏};
\node[block=gray, minimum width=92pt, minimum height=34pt] (reply) at (680,122) {\textbf{最终回复}};
\draw[flow] (user.east) -- (in.west);
\draw[flow] (in.east) -- (mid.west);
\draw[flow] (mid.east) -- (out.west);
\draw[flow] (out.east) -- (reply.west);
\node[block=orange, minimum width=150pt, minimum height=30pt] (ctx) at (388,62) {\textbf{统一风险上下文}：样本、轮次、标签、动作历史};
\draw[flow] (in.south) -- (ctx.north west);
\draw[flow] (mid.south) -- (ctx.north);
\draw[flow] (out.south) -- (ctx.north east);
\node[anchor=west, font=\bfseries\small] at (30,224) {三层防御总体架构图};
\end{tikzpicture}
}
\caption{三层防御总体架构图}
\label{fig:chapter5-defense-architecture}
\end{figure}

\subsection{防御决策与风险传递机制}

从实际实现方式看，本文分层防御并非由三个孤立模块简单串联，而是通过统一的风险上下文进行协同。系统在整个调用过程中维护统一的防御上下文，其中至少包含样本编号、模型名称、攻击类型、类别标签、轮次信息、原始提示、风险分数、风险标记、决策历史以及当前状态等字段。各层模块在处理时均以该上下文为输入，并在完成本层判断后更新共享的风险状态，从而形成贯穿全过程的风险传播链路。

在决策层面，系统支持多种防御动作，包括 \texttt{allow}、\texttt{rewrite}、\texttt{block}、\texttt{truncate}、\texttt{redact} 与 \texttt{replace} 等。不同动作对应不同强度的干预方式，其中允许表示不干预，改写表示保留调用但重写输入，阻断表示直接拒绝模型调用，截断表示限制交互继续推进，脱敏表示对输出中的危险细节进行局部遮蔽，替换表示以统一安全回复取代原始输出。为保证多层决策的一致性，系统采用“风险等级优先、动作强度次优先”的策略聚合原则：当不同层给出不同结论时，优先保留风险等级更高、动作更强的防御决策。

这一风险传递机制使本文方案具备两个特点。其一，防御不是单次静态判断，而是伴随模型调用全过程持续更新的动态过程；其二，前一层识别出的风险信号不会在进入下一层后丢失，而是被后续层继续利用。例如，输入层识别到的角色操控或规则豁免信号，可在交互层被解释为诱导风险，在输出层又进一步提升危险内容处理阈值。由此，三层防御得以形成“统一上下文、逐层增强、强动作优先”的整体设计。

本文将跨层风险聚合定义为：
\begin{equation}
R=\max \left(R_{\mathrm{in}}, R_{\mathrm{int}}, R_{\mathrm{out}}\right)
\end{equation}
其中，$R_{\mathrm{in}}$、$R_{\mathrm{int}}$ 与 $R_{\mathrm{out}}$ 分别表示输入层、交互层与输出层给出的风险分数。采用最大值聚合的原因在于，当任一层已识别出显著高风险信号时，系统整体不应因其他层判断较低而削弱干预强度。

在动作选择上，可进一步定义：
\begin{equation}
a^{*}=\arg\max_{a \in \mathcal{A}} \, \mathrm{priority}(a \mid R)
\end{equation}
其中，$\mathcal{A}=\{\texttt{allow},\texttt{rewrite},\texttt{block},\texttt{truncate},\texttt{redact},\texttt{replace}\}$ 表示动作集合，$\mathrm{priority}(a \mid R)$ 表示在当前风险水平下各动作的优先级。该设计对应“高风险优先、强动作次优先”的总体原则。

\begin{figure}[htbp]
\centering
\resizebox{0.98\textwidth}{!}{\begin{tikzpicture}[
    x=1pt,y=1pt,font=\small,>=stealth,
    block/.style={draw=#1!75!black, fill=#1!12, rounded corners=9pt, line width=0.9pt, align=center},
    flow/.style={->, line width=0.95pt, draw=gray!75}
]
\node[block=blue, minimum width=110pt, minimum height=38pt] (in) at (110,108) {\textbf{输入层}\\$R_{\mathrm{in}}$\\allow / rewrite / block};
\node[block=teal, minimum width=110pt, minimum height=38pt] (mid) at (296,108) {\textbf{交互层}\\$R_{\mathrm{int}}$\\warning / truncate};
\node[block=red, minimum width=110pt, minimum height=38pt] (out) at (482,108) {\textbf{输出层}\\$R_{\mathrm{out}}$\\redact / replace};
\node[block=orange, minimum width=116pt, minimum height=42pt] (agg) at (668,108) {\textbf{聚合决策}\\高风险优先\\强动作优先};
\draw[flow] (in.east) -- (mid.west);
\draw[flow] (mid.east) -- (out.west);
\draw[flow] (out.east) -- (agg.west);
\node[anchor=west, font=\scriptsize] at (52,48) {共享字段：risk score、risk flags、decision history、state、round index};
\node[anchor=west, font=\bfseries\small] at (18,170) {风险传递与强动作优先机制图};
\end{tikzpicture}
}
\caption{风险传递与强动作优先机制图}
\label{fig:chapter5-risk-propagation}
\end{figure}

\section{输入层防御机制}

\subsection{输入层防御的功能定位}

输入层防御的主要任务是在模型推理发生前完成风险识别与前置干预，其核心目标是尽量在攻击尚未进入生成阶段之前切断明显恶意请求。结合第四章实验结果，输入层主要针对以下几类高风险攻击信号：显式危险意图、提示注入与规则劫持、角色操控与人格伪装、规则豁免与越权表述，以及与违法活动、网络入侵、金融犯罪和其他高危任务直接相关的关键词组合。

从攻击面映射上看，输入层尤其对应第四章中表现突出的认知误导、角色扮演、对抗后缀、结构伪装及部分指令重构类攻击。这些攻击虽然实现形式不同，但其中大量样本都在提示文本本身保留了可识别的风险信号，例如“忽略之前的规则”“你现在扮演某个角色”“不必遵守限制”“请自由表达”“如何设立非法业务”等。因此，输入层的研究重点不是判断模型已经输出了什么，而是判断当前输入是否已经具有明显越狱倾向。

\begin{figure}[htbp]
\centering
\resizebox{0.98\textwidth}{!}{\begin{tikzpicture}[
    x=1pt,y=1pt,font=\small,>=stealth,
    block/.style={draw=#1!75!black, fill=#1!12, rounded corners=9pt, line width=0.9pt, align=center},
    flow/.style={->, line width=0.95pt, draw=gray!75}
]
\node[block=gray, minimum width=90pt, minimum height=34pt] (recv) at (82,90) {\textbf{输入接收}};
\node[block=blue, minimum width=96pt, minimum height=34pt] (rule) at (212,90) {\textbf{规则匹配}};
\node[block=teal, minimum width=96pt, minimum height=34pt] (cls) at (342,90) {\textbf{分类器打分}};
\node[block=orange, minimum width=96pt, minimum height=34pt] (fuse) at (472,90) {\textbf{风险融合}};
\node[block=red, minimum width=110pt, minimum height=42pt] (act) at (624,90) {\textbf{动作输出}\\allow / rewrite / block};
\draw[flow] (recv.east) -- (rule.west);
\draw[flow] (rule.east) -- (cls.west);
\draw[flow] (cls.east) -- (fuse.west);
\draw[flow] (fuse.east) -- (act.west);
\node[anchor=west, font=\bfseries\small] at (18,146) {输入层防御流程图};
\node[anchor=west, font=\scriptsize, text=gray!70!black] at (18,28) {映射重点：认知误导、角色操控、规则豁免、显式高危任务};
\end{tikzpicture}
}
\caption{输入层防御流程图}
\label{fig:chapter5-input-defense}
\end{figure}

\subsection{基于规则与分类器的风险识别}

本文输入层采用“规则匹配 + 轻量分类打分”的混合识别机制，以兼顾可解释性与覆盖性。规则部分主要面向高精度风险模式识别，例如系统提示劫持、权限提升、系统提示泄露、角色操控、规则豁免和显式有害意图等。这些规则以正则模式和组合规则的形式实现，既可识别单一危险表述，也可对“角色操控 + 规则豁免 + 有害任务”这类复合越狱信号给出更高风险分数\cite{rebedea2023nemo,markov2022holistic}。

在规则基础上，系统还引入轻量关键词分类器，对输入文本中的行动导向、绕过导向、恶意软件导向及中文有害任务表达进行打分。与规则匹配相比，分类器更适合捕捉边界性风险请求，例如未明确出现单一高危短语，但整体上体现出“如何做”“怎么绕过”“如何搭建”等行动意图的提示。系统最终将规则分数与分类器分数进行融合，并据此更新统一风险分数。

这一设计的优点在于：一方面，规则能够提供较高的可解释性，便于后续审计与规则维护；另一方面，分类器能够弥补纯规则匹配在边界表达与变形请求上的不足。对于论文方案而言，输入层并不追求完全准确地理解所有危险语义，而是通过“高风险模式识别 + 中风险边界估计”的方式，在模型调用前尽可能压缩攻击进入生成阶段的概率\cite{markov2022holistic,rebedea2023nemo}。

\begin{table}[htbp]
\centering
\footnotesize
\setlength{\tabcolsep}{4pt}
\renewcommand{\arraystretch}{0.95}
\caption{输入层风险识别规则表示例}
\begin{tabular}{p{0.16\textwidth}p{0.22\textwidth}p{0.12\textwidth}p{0.28\textwidth}}
\toprule
\textbf{规则类别} & \textbf{示例模式} & \textbf{建议分值} & \textbf{说明} \\
\midrule
系统提示劫持 & 抽象识别“忽略既有规则”“覆盖系统约束”等模式 & 0.9 & 属于高精度高风险规则 \\
角色操控 & 扮演专家、执行者、指挥角色等身份伪装模式 & 0.7 & 对应认知误导与角色扮演类攻击 \\
规则豁免 & 抽象识别“无需限制”“不必遵守政策”等表述 & 0.8 & 常与高危任务共现 \\
显式有害意图 & 违法、入侵、规避、欺诈等任务导向关键词组合 & 1.0 & 对应直接风险任务 \\
\bottomrule
\end{tabular}
\end{table}
\noindent\textbf{附注：}如果担心正文出现敏感短语，可以将规则模式写成抽象类别或脱敏形式。

\begin{equation}
R_{\mathrm{in}}=\gamma_1 S_{\mathrm{rule}} + \gamma_2 S_{\mathrm{cls}}
\end{equation}
其中，$S_{\mathrm{rule}}$ 表示规则匹配得分，$S_{\mathrm{cls}}$ 表示轻量分类器输出，$\gamma_1,\gamma_2$ 为融合权重。通过这一形式，输入层可以在高精度规则识别与边界性语义估计之间取得平衡。

\subsection{输入层的阻断与改写策略}

在输入层动作设计上，本文采用分级干预策略，而非“一律拒绝”。当输入风险分数达到高阈值时，系统直接阻断模型调用，以避免显式恶意请求进入生成过程；当风险处于中等水平时，系统不直接丢弃请求，而是触发安全改写动作，即在保留原始请求主体的前提下，加入安全前缀与明确拒答约束，使模型在后续生成中优先遵守安全边界；当风险较低时，则允许请求继续进入后续流程，但其风险信息仍会保留在统一上下文中，以供交互层和输出层参考。

这种设计兼顾了安全性与可用性。对于明确高危输入，直接阻断可以最大限度降低风险；对于边界性或语义不够清晰的输入，安全改写则可以减少误伤正常请求的概率。尤其考虑到第四章中大量样本表现为“不确定输出”，输入层改写机制实际上承担了一个重要角色，即在调用前主动加强安全语境，减少模型进入模糊边界状态的可能性\cite{rebedea2023nemo}。

\begin{table}[htbp]
\centering
\footnotesize
\setlength{\tabcolsep}{4pt}
\renewcommand{\arraystretch}{0.95}
\caption{输入层动作决策表}
\begin{tabular}{p{0.16\textwidth}p{0.18\textwidth}p{0.18\textwidth}p{0.28\textwidth}}
\toprule
\textbf{风险区间} & \textbf{动作} & \textbf{是否继续调用模型} & \textbf{说明} \\
\midrule
低风险 & \texttt{allow} & 是 & 记录风险上下文，进入后续流程 \\
中风险 & \texttt{rewrite} & 是 & 加入安全前缀或改写提示 \\
高风险 & \texttt{block} & 否 & 直接阻断请求 \\
\bottomrule
\end{tabular}
\end{table}

\section{交互层防御机制}

\subsection{交互层防御的功能定位}

交互层防御主要面向过程性攻击风险，其核心任务是监测风险如何在对话轮次或连续上下文中累积，并在必要时对交互链路进行限制。即使单轮输入本身风险不高，攻击者仍可能通过连续试探、角色维持、规则豁免强化和多轮追问逐步扩大模型的偏移程度。因此，交互层并不是对输入层的重复，而是针对“攻击正在形成”的过程实施动态控制。

结合本文前述实验设计和未来扩展方向，交互层重点对应上下文操控、多轮越狱以及静态单轮中通过长前缀、问答链和任务序列实现的结构诱导风险。其研究意义在于，安全对齐的失效并不总是发生在单个输入点上，而可能表现为风险信号在多轮中逐步聚集、在上下文中持续强化，并最终在某一轮转化为明确突破。

\begin{figure}[htbp]
\centering
\resizebox{0.98\textwidth}{!}{\begin{tikzpicture}[
    x=1pt,y=1pt,font=\small,>=stealth,
    state/.style={draw=#1!75!black, fill=#1!12, rounded corners=18pt, line width=1pt, align=center, minimum width=92pt, minimum height=42pt},
    flow/.style={->, line width=0.95pt, draw=gray!75}
]
\node[state=blue] (s1) at (96,92) {\textbf{正常}};
\node[state=orange] (s2) at (262,92) {\textbf{预警}};
\node[state=teal] (s3) at (428,92) {\textbf{受限}};
\node[state=red] (s4) at (594,92) {\textbf{阻断}};
\draw[flow] (s1.east) -- node[above, font=\scriptsize] {$R_t>\tau_1$} (s2.west);
\draw[flow] (s2.east) -- node[above, font=\scriptsize] {$R_t>\tau_2$} (s3.west);
\draw[flow] (s3.east) -- node[above, font=\scriptsize] {$R_t>\tau_3$} (s4.west);
\draw[flow] (s2.south) .. controls (262,36) and (96,36) .. node[below, font=\scriptsize] {风险回落} (s1.south);
\node[anchor=west, font=\bfseries\small] at (18,154) {交互层过程控制图};
\end{tikzpicture}
}
\caption{交互层过程控制图}
\label{fig:chapter5-interaction-control}
\end{figure}

\subsection{风险累积与诱导行为检测}

本文交互层的核心机制是基于统一风险分数、轮次信息与诱导性标记进行动态判断。系统将输入层识别出的角色操控、规则豁免、复合越狱模式等风险标记保留在上下文中，并结合当前轮次与累计风险分数判断攻击是否已进入诱导阶段。当风险分数超过警戒阈值时，系统将当前交互状态标记为预警；若在此基础上继续出现多轮诱导信号，或风险分数进一步上升至更高区间，则系统会将交互视为受限状态甚至直接阻断\cite{chao2024pair,anil2024manyshot}。

从方案设计角度看，这一机制体现了两个重要原则。第一，风险不是静态的，而是可积累的。一次中风险输入未必足以触发强动作，但若相同类型的诱导信号在后续轮次中持续出现，则其总体风险应显著提高。第二，轮次本身也是一种风险因素。当交互进入更后期轮次时，系统对风险的容忍度应适当降低，因为攻击者在真实场景中往往利用多轮试探逐步突破边界。

\begin{equation}
R_t=\alpha R_{t-1}+\delta_t+\eta t
\end{equation}
其中，$\delta_t$ 表示当前轮风险增量，$\eta$ 表示轮次惩罚系数，$t$ 表示当前轮次。该定义体现了交互层“风险可累积、轮次本身也是风险因素”的基本思想。

\begin{equation}
\mathrm{state}_t=
\begin{cases}
\texttt{normal}, & R_t < \tau_1 \\
\texttt{warning}, & \tau_1 \le R_t < \tau_2 \\
\texttt{restricted}, & \tau_2 \le R_t < \tau_3 \\
\texttt{blocked}, & R_t \ge \tau_3
\end{cases}
\end{equation}
其中，$\tau_1,\tau_2,\tau_3$ 为逐级递增阈值，对应正常、预警、受限与阻断四种交互状态。

\subsection{多轮场景下的限制与截断策略}

在动作层面，交互层主要承担两类干预职责。其一是在风险达到警戒阈值但尚未严重失控时，保留模型调用但将系统状态标记为预警，从而为后续输出层提供更严格的上下文；其二是在检测到明显诱导性多轮行为或超出允许轮次的高风险交互时，直接对会话链进行截断，阻止攻击继续积累。对于极高风险状态，交互层也可直接返回阻断决策，从而在过程阶段提前终止对话。

这一设计使交互层成为输入层与输出层之间的重要缓冲带。它既避免了所有中风险对话都被立即封禁，也避免了攻击在连续轮次中无约束地发展。尤其在未来补充多轮越狱实验结果后，交互层将是验证“过程性防御是否有效”的关键模块，因此本章先从方案角度对其进行完整说明。

\begin{table}[htbp]
\centering
\footnotesize
\setlength{\tabcolsep}{4pt}
\renewcommand{\arraystretch}{0.95}
\caption{交互层动作触发表}
\begin{tabular}{p{0.15\textwidth}p{0.20\textwidth}p{0.22\textwidth}p{0.25\textwidth}}
\toprule
\textbf{交互状态} & \textbf{动作} & \textbf{适用场景} & \textbf{说明} \\
\midrule
正常 & \texttt{allow} & 单轮低风险或早期低风险对话 & 仅记录上下文 \\
预警 & \texttt{allow} + warning flag & 中风险多轮或持续诱导早期阶段 & 加强后续层敏感度 \\
受限 & \texttt{truncate} & 风险持续累计且接近突破边界 & 限制后续追问 \\
阻断 & \texttt{block} & 高风险或明显多轮诱导 & 提前终止会话 \\
\bottomrule
\end{tabular}
\end{table}

\section{输出层防御机制}

\subsection{输出层防御的功能定位}

输出层防御是整个分层体系中的最后一道安全保障，其任务是在模型已经给出回复之后，对回复文本中的危险内容进行再识别、再处理与审计记录。考虑到前两层防御无法保证在所有场景下都完全阻断风险，输出层的存在并不是对前置机制失效的被动补救，而是整个体系中不可或缺的终端兜底层\cite{markov2022holistic,phute2024llmselfdefense}。

从攻击面对应关系看，输出层覆盖所有可能穿透前置防线的攻击路径，但其重点关注的是高危可操作内容，例如攻击步骤、规避执法建议、网络入侵细节、金融犯罪操作、暴力犯罪流程、毒品交易流程与其他明显具有执行导向的危险表达。换言之，输出层不再以“用户想做什么”为主要判断对象，而转向“模型已经输出了什么”。

\begin{figure}[htbp]
\centering
\resizebox{0.98\textwidth}{!}{\begin{tikzpicture}[
    x=1pt,y=1pt,font=\small,>=stealth,
    block/.style={draw=#1!75!black, fill=#1!12, rounded corners=9pt, line width=0.9pt, align=center},
    flow/.style={->, line width=0.95pt, draw=gray!75}
]
\node[block=gray, minimum width=100pt, minimum height=34pt] (raw) at (84,92) {\textbf{模型原始回复}};
\node[block=blue, minimum width=110pt, minimum height=34pt] (detect) at (236,92) {\textbf{危险内容检测}};
\node[block=orange, minimum width=110pt, minimum height=34pt] (enhance) at (388,92) {\textbf{上下文增强判断}};
\node[block=red, minimum width=118pt, minimum height=42pt] (act) at (554,92) {\textbf{动作输出}\\allow / redact / replace};
\node[block=teal, minimum width=96pt, minimum height=34pt] (audit) at (690,92) {\textbf{审计写入}};
\draw[flow] (raw.east) -- (detect.west);
\draw[flow] (detect.east) -- (enhance.west);
\draw[flow] (enhance.east) -- (act.west);
\draw[flow] (act.east) -- (audit.west);
\node[anchor=west, font=\bfseries\small] at (18,154) {输出层内容安全兜底图};
\end{tikzpicture}
}
\caption{输出层内容安全兜底图}
\label{fig:chapter5-output-defense}
\end{figure}

\subsection{面向危险内容的输出风险识别}

本文输出层主要通过规则化内容识别实现风险判断。其检测重点并非仅限于单个关键词，而是围绕“步骤化、操作化、规避化、组织化”的危险输出特征展开。例如，系统会关注是否出现明确步骤标记、攻击链细节、会话凭证窃取、提权与持久化术语、金融犯罪资金路由、规避审计与监管的操作建议，以及毒品前体、藏匿与跨境转运等表达。这种设计使输出层更接近内容安全审查，而不是简单的词汇屏蔽\cite{markov2022holistic,phute2024llmselfdefense}。

此外，输出层并非孤立判断文本本身，而会结合前面阶段的上下文状态进行风险增强。若当前样本在输入层或交互层已被识别出角色操控、规则豁免或高风险组合信号，则输出层会进一步提高对危险内容的敏感性。这意味着同样一段边界内容，在普通对话与越狱上下文中将被赋予不同风险权重，从而提升输出层对复杂攻击场景的适应能力\cite{rebedea2023nemo,phute2024llmselfdefense}。

\begin{equation}
R_{\mathrm{out}}=S_{\mathrm{content}}+\rho \cdot C_{\mathrm{context}}
\end{equation}
其中，$S_{\mathrm{content}}$ 表示输出内容审查得分，$C_{\mathrm{context}}$ 表示来自输入层与交互层的上下文风险，$\rho$ 为增强系数。该式说明输出层并非只看文本本身，而是结合前置阶段风险信号进行增强判断。

\begin{equation}
a_{\mathrm{out}}=
\begin{cases}
\texttt{allow}, & R_{\mathrm{out}} < \lambda_1 \\
\texttt{redact}, & \lambda_1 \le R_{\mathrm{out}} < \lambda_2 \\
\texttt{replace}, & R_{\mathrm{out}} \ge \lambda_2
\end{cases}
\end{equation}
其中，$\lambda_1,\lambda_2$ 为分级处理阈值，分别对应放行、脱敏与替换三档动作。

\subsection{输出替换、脱敏与审计记录策略}

在动作设计上，输出层采用“高风险替换，中风险脱敏，低风险放行”的分级处理策略。对于高风险输出，系统直接以统一安全回复替换原始内容，避免高危信息暴露给最终用户；对于中风险输出，则对其中的敏感术语、步骤性表达或操作细节进行局部脱敏，在尽可能保留低风险信息的同时消除可操作危险内容；对于低风险输出，则允许其原样返回。

与输入层和交互层相比，输出层还额外承担审计功能。系统会将模型名称、样本编号、防御动作、风险等级、触发原因、原始回复与处理后回复等信息写入审计存储，以便后续进行误报复核、规则优化和防御效果分析。对于论文方案而言，这一设计的重要性在于：防御不仅要“拦住”，还要“可解释、可追踪、可改进”\cite{markov2022holistic,rebedea2023nemo}。

\begin{table}[htbp]
\centering
\footnotesize
\setlength{\tabcolsep}{4pt}
\renewcommand{\arraystretch}{0.95}
\caption{输出层审计字段设计表}
\begin{tabular}{p{0.20\textwidth}p{0.18\textwidth}p{0.42\textwidth}}
\toprule
\textbf{字段名} & \textbf{类型} & \textbf{说明} \\
\midrule
model\_name & string & 模型名称 \\
sample\_id & string/int & 样本编号 \\
action & string & 最终防御动作 \\
risk\_level & int & 输出风险等级 \\
trigger\_reason & string & 触发规则或风险原因摘要 \\
raw\_response & text & 原始回复文本 \\
processed\_response & text & 处理后回复文本 \\
timestamp & datetime & 记录时间 \\
\bottomrule
\end{tabular}
\end{table}

\section{三层防御的协同作用与适用范围}

\subsection{分层协同的必要性分析}

从整体上看，输入层、交互层与输出层分别作用于攻击链的不同阶段，三者之间并不存在功能替代关系，而是构成了互补的协同体系。输入层负责识别显式越狱信号，适合处理攻击进入前的高危样本；交互层负责监测状态累积和过程性诱导，适合处理连续对话和上下文扩散；输出层负责兜底危险内容，适合处理前置阶段未能完全覆盖的残余风险。三层联动的价值在于，即便某一层未能完全拦截风险，其他层仍可继续发挥补偿作用，从而整体提高系统鲁棒性。

结合第四章结果，这种分层设计具有明确必要性。由于静态实验中大量样本表现为不确定状态，说明风险边界并非仅存在“是否进入”这一单一问题，而是贯穿输入识别、上下文推进与输出内容三个环节。单一规则过滤难以同时覆盖这些环节，必须通过多层机制分别应对不同类型的脆弱性来源。

\subsection{面向多维攻击框架的防御映射关系}

从攻击面映射角度看，本文分层防御方案与前文多维越狱攻击框架保持了较强的一致性。对于角色扮演、认知误导、对抗后缀、结构伪装和显式危险意图等攻击，输入层承担主要识别职责；对于指令重构、上下文前缀、多轮试探与连续诱导等攻击，交互层承担主要约束职责；而无论攻击来源如何，一旦输出中出现步骤化、操作化或规避性危险内容，输出层都将作为最终兜底机制介入。

这一映射关系说明，本文提出的分层防御并非脱离实验结果而抽象设计，而是直接回应第三章与第四章中所揭示的关键脆弱维度。其本质上是将“攻击作用于何处”的分析结果转化为“防御在哪一层介入”的系统设计，从而实现攻防结构上的对应。

\begin{figure}[htbp]
\centering
\resizebox{0.98\textwidth}{!}{\begin{tikzpicture}[
    x=1pt,y=1pt,font=\small,>=stealth,
    block/.style={draw=#1!75!black, fill=#1!12, rounded corners=9pt, line width=0.9pt, align=center},
    flow/.style={->, line width=0.9pt, draw=gray!70}
]
\node[block=gray, minimum width=118pt, minimum height=34pt] (a1) at (110,170) {\textbf{认知误导 / 角色操控}};
\node[block=gray, minimum width=118pt, minimum height=34pt] (a2) at (110,118) {\textbf{指令重构 / 上下文前缀}};
\node[block=gray, minimum width=118pt, minimum height=34pt] (a3) at (110,66) {\textbf{表面扰动 / 编码变换}};

\node[block=blue, minimum width=106pt, minimum height=34pt] (d1) at (360,170) {\textbf{输入层}};
\node[block=teal, minimum width=106pt, minimum height=34pt] (d2) at (360,118) {\textbf{交互层}};
\node[block=red, minimum width=106pt, minimum height=34pt] (d3) at (360,66) {\textbf{输出层}};

\node[block=orange, minimum width=128pt, minimum height=42pt] (goal) at (604,118) {\textbf{防御目标}\\早期阻断 / 过程限制 / 输出兜底};

\draw[flow] (a1.east) -- (d1.west);
\draw[flow] (a1.east) .. controls (210,170) and (286,118) .. (d2.west);
\draw[flow] (a2.east) -- (d2.west);
\draw[flow] (a2.east) .. controls (220,118) and (286,66) .. (d3.west);
\draw[flow] (a3.east) -- (d1.west);
\draw[flow] (a3.east) .. controls (230,66) and (286,66) .. (d3.west);
\draw[flow] (d1.east) -- (goal.west);
\draw[flow] (d2.east) -- (goal.west);
\draw[flow] (d3.east) -- (goal.west);
\node[anchor=west, font=\bfseries\small] at (18,220) {攻击维度与防御层映射图};
\end{tikzpicture}
}
\caption{攻击维度与防御层映射图}
\label{fig:chapter5-attack-defense-map}
\end{figure}

\section{防御效果评估方案}

\subsection{防御实验设置}

由于当前防御状态下的完整实验结果尚未生成，本节暂不报告具体防御效果，仅说明后续实验应如何组织。为完成防御验证，至少需要在与第四章一致的评测口径下补充以下数据：第一，启用防御前后在同一静态测试集上的完整对照结果，包括三模型在九类静态攻击维度下的统一 \texttt{records.csv} 与 \texttt{group\_metrics.csv}；第二，自动化红队样本在启用防御后的对照结果，以评估输入层与输出层对动态变体的适应能力；第三，多轮交互场景下的统一轮次统计文件，以验证交互层对首次成功轮次与累计成功率的抑制效果。

在实验组织上，建议至少设置“无防御”“仅输入层”“输入层 + 交互层”“三层全开”四种对照条件，以区分各层单独作用与全层协同作用的贡献。同时，所有实验均应复用第三章与第四章中已建立的统一数据与判定流程，避免由于评测口径变化而影响结果可比性。

\begin{table}[htbp]
\centering
\footnotesize
\setlength{\tabcolsep}{4pt}
\renewcommand{\arraystretch}{0.95}
\caption{防御对照实验设置表}
\begin{tabular}{p{0.18\textwidth}p{0.14\textwidth}p{0.14\textwidth}p{0.14\textwidth}p{0.24\textwidth}}
\toprule
\textbf{实验条件} & \textbf{输入层} & \textbf{交互层} & \textbf{输出层} & \textbf{说明} \\
\midrule
无防御 & 关闭 & 关闭 & 关闭 & 基线条件 \\
仅输入层 & 开启 & 关闭 & 关闭 & 评估前置拦截能力 \\
输入层 + 交互层 & 开启 & 开启 & 关闭 & 评估过程控制能力 \\
三层全开 & 开启 & 开启 & 开启 & 评估完整协同防御效果 \\
\bottomrule
\end{tabular}
\end{table}

\subsection{防御效果评估指标}

在防御效果评估中，建议重点关注以下指标：第一，防御后越狱成功率的下降幅度，用于衡量系统总体拦截能力；第二，高风险输出占比的变化，用于评估输出层兜底对危险内容强度的抑制效果；第三，不确定样本比例的变化，用于观察防御是否能够将边界模糊状态收敛为更明确的安全拒答；第四，不同防御动作的触发频率与分布，用于分析输入层、交互层与输出层的实际介入路径；第五，误拦截率与改写率，用于衡量防御对正常请求可用性的影响；第六，多轮场景下首次成功轮次的延后程度与累计成功率的压缩程度，用于评估交互层的过程性控制能力。

这些指标的共同目标，不仅是回答“拦住了多少”，也要回答“以什么代价拦住了多少”。若防御机制仅通过大规模阻断实现低成功率，则其工程价值有限；相反，若系统能够在有限误伤代价下显著压缩高脆弱攻击面的突破空间，则更能体现分层防御方案的实际意义。

\begin{figure}[htbp]
\centering
\resizebox{0.98\textwidth}{!}{\begin{tikzpicture}[
    x=1pt,y=1pt,font=\small,>=stealth,
    block/.style={draw=#1!75!black, fill=#1!12, rounded corners=9pt, line width=0.9pt, align=center},
    flow/.style={->, line width=0.95pt, draw=gray!75}
]
\node[block=blue, minimum width=108pt, minimum height=38pt] (base) at (96,98) {\textbf{无防御基线}};
\node[block=orange, minimum width=108pt, minimum height=38pt] (single) at (258,98) {\textbf{单层防御对照}};
\node[block=teal, minimum width=108pt, minimum height=38pt] (full) at (420,98) {\textbf{三层协同防御}};
\node[block=red, minimum width=128pt, minimum height=48pt] (metrics) at (618,98) {\textbf{评估指标}\\ASR / UR / HR\\误报率 / 可用性损失};
\draw[flow] (base.east) -- (single.west);
\draw[flow] (single.east) -- (full.west);
\draw[flow] (full.east) -- (metrics.west);
\node[anchor=west, font=\bfseries\small] at (18,162) {防御效果评估框架图};
\node[anchor=west, font=\scriptsize, text=gray!70!black] at (18,30) {用于替代未完成的结果图组, 明确后续对照实验的指标与比较路径};
\end{tikzpicture}
}
\caption{防御效果评估框架图}
\label{fig:chapter5-evaluation-framework}
\end{figure}

\subsection{结果分析位置预留}

待防御实验结果补齐后，本节建议补入以下内容：第一，启用防御前后的分维度成功率对照表与高风险占比对照表；第二，输入层、交互层和输出层各自的动作分布图，展示不同攻击维度主要被哪一层拦截；第三，多轮场景下累计成功率曲线与首次成功轮次分布图，用于展示交互层对过程性攻击的抑制效果；第四，典型防御案例分析，用于说明输入改写、会话截断与输出替换在具体场景中的作用机制。

在相关数据尚未齐备之前，本文仅将这些内容作为后续实验填充位置予以保留，不对防御效果做超出证据范围的结论性陈述。

\begin{table}[htbp]
\centering
\footnotesize
\setlength{\tabcolsep}{4pt}
\renewcommand{\arraystretch}{0.95}
\caption{典型防御案例分析表}
\begin{tabular}{p{0.10\textwidth}p{0.18\textwidth}p{0.16\textwidth}p{0.16\textwidth}p{0.24\textwidth}}
\toprule
\textbf{案例编号} & \textbf{攻击类型} & \textbf{触发层级} & \textbf{最终动作} & \textbf{分析重点} \\
\midrule
Def-1 & 认知误导 / 角色操控类输入 & 输入层 & \texttt{rewrite} 或 \texttt{block} & 输入改写如何抬高安全边界 \\
Def-2 & 上下文诱导 / 多轮试探类攻击 & 交互层 & \texttt{truncate} 或 \texttt{block} & 会话截断如何抑制过程性诱导 \\
Def-3 & 已生成边界性危险回复 & 输出层 & \texttt{redact} 或 \texttt{replace} & 输出替换如何消除高危操作细节 \\
\bottomrule
\end{tabular}
\end{table}

\section{本章小结}

本章面向越狱攻击的全过程风险控制问题，提出了由输入层、交互层与输出层构成的三层防御方案。该方案以统一风险上下文为基础，通过前置识别、过程约束与输出兜底三个阶段实现协同防护，既能够对应第四章中暴露出的认知误导、指令重构、上下文操控等高脆弱攻击面，也为后续自动化红队与多轮交互场景预留了兼容的防御接口。

需要强调的是，当前章节的重点在于完整说明防御方案的结构设计与机制逻辑，而非提前宣称其实验效果。由于防御状态下的系统评测结果尚未全部生成，相关效果分析仍需在后续补充实验后进一步完成。尽管如此，现有方案已经在结构上形成了与前文攻击框架相对应的分层防护思路，为全文从“漏洞分析—攻击评测—防御设计”构建闭环研究路径提供了关键支撑。
